\documentclass[11pt]{article}
\usepackage[a4paper,margin=1in]{geometry}
\usepackage{amsmath,amssymb,mathtools,bm}
\usepackage{enumitem,booktabs,array}
\usepackage{tcolorbox}
\usepackage{hyperref}
\usepackage[protrusion=true,expansion=false]{microtype}
\hypersetup{colorlinks=true,linkcolor=blue,urlcolor=blue,citecolor=blue}
\setlist[itemize]{topsep=2pt,itemsep=2pt}
\setlist[enumerate]{topsep=2pt,itemsep=2pt}
\newtcolorbox{keybox}{colback=blue!3!white,colframe=blue!40!black,boxrule=0.4pt,arc=1mm}
\newtcolorbox{alertbox}{colback=red!3!white,colframe=red!40!black,boxrule=0.4pt,arc=1mm}
\newtcolorbox{answerbox}{colback=green!3!white,colframe=green!40!black,boxrule=0.4pt,arc=1mm}
\newtcolorbox{drillbox}{colback=yellow!5!white,colframe=orange!60!black,boxrule=0.4pt,arc=1mm}
\newcommand{\EV}{\mathbb{E}}
\newcommand{\Var}{\mathrm{Var}}
\newcommand{\Cov}{\mathrm{Cov}}
\newcommand{\blank}[1]{\underline{\hspace{#1}}}

\title{W13-01: Murfin \& Spiegel (2020, RFS) \\
\large ``Is the Risk of Sea Level Rise Capitalized in Residential Real Estate?'' --- Active Recall Workbook}
\author{Banking \& Risk Management --- Exam Prep}
\date{\today}
\begin{document}\maketitle

\begin{keybox}
\textbf{Paper in one sentence.} Once one controls for local ground subsidence (which is highly correlated with SLR exposure), the Bernstein-Gustafson-Lewis $\sim$7\% price discount on SLR-exposed coastal homes falls essentially to zero --- MS argue the housing market does \emph{not} meaningfully price long-horizon SLR risk in the US cross-section they examine.

\textbf{Placement.} The canonical critical counterpoint to Bernstein-Gustafson-Lewis (W12-03). Keys-Mulder (W14-03) provides updated evidence that markets are slowly adjusting as information improves.
\end{keybox}

\section*{Round 1: Complete Walk-Through}

\subsection*{1.1 Critical idea}
Some coastal areas are exposed to SLR partly because they are physically \emph{sinking} (subsidence) --- e.g., Louisiana, parts of Florida, Virginia. Subsidence is itself a strong predictor of house prices (structural damage, flooding) and is correlated with NOAA SLR inundation maps. Omitting subsidence could upward-bias the SLR discount.

\subsection*{1.2 Data and design}
Same housing-transaction data as BGL (ZIP-level), 2007--2016, combined with US Geological Survey subsidence estimates at fine spatial resolution. Augmented hedonic regression:
\[
\ln P_i=\alpha_{z,t}+\beta\cdot\text{SLR}_i+\delta\cdot\text{Subsidence}_i+X_i'\gamma+\varepsilon_i.
\]

\subsection*{1.3 Headline results}
\begin{itemize}
\item Without subsidence control: SLR discount $\sim$$-7\%$ (matching BGL).
\item With subsidence control: SLR discount falls to $\sim$$-1$ to $-2\%$, often statistically indistinguishable from zero.
\item Subsidence itself carries a large, significant negative coefficient.
\item Interpretation: the ``SLR discount'' in BGL largely reflects subsidence-driven present damage, not future climate risk.
\end{itemize}

\subsection*{1.4 Implication}
\begin{itemize}
\item Long-horizon SLR risk is under-priced in the 2007--2016 cross-section.
\item Consistent with MS's earlier work on under-pricing of long-horizon risks generally.
\item Suggests belief / time-horizon / information frictions in climate-risk pricing.
\end{itemize}

\subsection*{1.5 Caveats}
\begin{alertbox}
\begin{itemize}
\item Subsidence and SLR projections have physical commonalities; decomposition is imperfect.
\item Sample ends 2016; later data (Keys-Mulder 2025) show growing discounts.
\item Bayesian interpretation: results depend on how much weight one puts on long-horizon versus present risk.
\end{itemize}
\end{alertbox}

\section*{Round 2: Level 1 Fill-in}
\begin{enumerate}
\item Omitted variable flagged by MS: \blank{2cm}.
\item Baseline (BGL-style) discount: $\sim$\blank{1cm}\%.
\item With subsidence control: $\sim$\blank{1cm}--\blank{1cm}\%.
\item Interpretation: SLR risk is \blank{1cm}-priced.
\item Updated evidence in Keys-\blank{1.5cm} (2025).
\end{enumerate}

\section*{Round 3: Level 2 Prompts}
\begin{enumerate}
\item Write the augmented hedonic and identify the potential bias in BGL.
\item Why is subsidence so strongly correlated with SLR exposure?
\item Discuss implications for welfare calculations of climate adaptation.
\item How do Bernstein-Billings-Gustafson-Lewis (2022) reconcile their discount with MS's null?
\end{enumerate}

\section*{Round 4: Minimal Scaffolding}
From a blank page: (i) omitted-variable argument, (ii) spec, (iii) magnitude comparison, (iv) interpretation, (v) caveats.

\section*{Round 5: Blank Page}
Essay: is climate risk priced in US residential real estate? BGL vs.\ MS.

\vspace{6pt}
\begin{drillbox}
\textbf{Speed Drill.} (1) Omitted variable. (2) Baseline discount. (3) Controlled discount. (4) Implication for pricing. (5) Comparison paper.
\end{drillbox}

\section*{Quick Reference \& Answer Key}
\begin{answerbox}
\begin{itemize}
\item \textbf{Omitted variable.} Ground subsidence.
\item \textbf{Baseline.} $\sim$$-7\%$ (BGL 2019).
\item \textbf{Controlled.} $\sim$$-1$ to $-2\%$, often $\approx 0$.
\item \textbf{Read.} SLR risk is under-priced in 2007--16 data.
\end{itemize}
\end{answerbox}

\begin{keybox}
\textbf{30-second exam pitch.} Murfin \& Spiegel (2020) re-examine the Bernstein-Gustafson-Lewis $\sim$7\% sea-level-rise discount in coastal home prices and find that once one controls for local ground subsidence (using USGS data), the SLR discount shrinks to $-1$ to $-2\%$, often statistically indistinguishable from zero. The insight is that subsidence correlates strongly with SLR-inundation classification: the BGL discount largely reflects present subsidence-driven damage, not forward-looking SLR expectations. MS argue that long-horizon climate risk is under-priced in the 2007--2016 US cross-section, consistent with belief, information, and time-horizon frictions. The paper is the key critical counterpoint to BGL; Bernstein-Billings-Gustafson-Lewis (2022) reconciles by emphasizing that belief heterogeneity, not just physical risk, drives pricing, and Keys-Mulder (2025) shows more recent data indicate slow catch-up as information improves.
\end{keybox}

\end{document}
